% Chapter 4: GPIO Implementation
% File: chapters/chapter-04-gpio/chapter-04-gpio.tex

\chapter{Phase A: GPIO Implementation}
\label{ch:chapter-04}
\chapthumb{4}

\begin{quote}
    \textit{``A pin is not an abstraction problem. It is a register problem that everybody thinks is
    an abstraction problem.''}
\end{quote}

A GPIO is the simplest peripheral on any microcontroller and the one that produces the most
over-engineered code. This chapter takes the opposite route: it starts from the registers on this
part, derives the smallest abstraction that is still safe, and shows the three defects that the
bigger abstraction — 605 lines of type-state pins — did not prevent.

Everything here is measured on the XIAO MG24 Sense (EFR32MG24B220F1536IM48-B), and the numbers are
read back over SWD rather than assumed.

\minitoc

\clearpage

\section{The Registers, and How Few There Are}
\label{sec:gpio-registers}

A GPIO port on the EFR32MG24 is not a black box. It is five things:

\begin{table}[htbp]
    \centering
    \footnotesize
    \begin{tabular}{@{}p{0.18\textwidth}p{0.42\textwidth}p{0.30\textwidth}@{}}
        \toprule
        \textbf{Register} & \textbf{What it decides} & \textbf{Access} \\
        \midrule
        \codeinline{MODEL} & mode of pins 0--7 & read-modify-write, per pin \\
        \codeinline{MODEH} & mode of pins 8--15 & only on PORTA and PORTC here \\
        \codeinline{DOUT} & the output latch, one bit per pin & read-modify-write \\
        \codeinline{DIN} & the pin as the outside world holds it & read only \\
        \codeinline{CMU\_CLKEN0} & whether the port answers at all & gate before the first access \\
        \bottomrule
    \end{tabular}
    \caption{The GPIO port, in full. Five registers, and only one of them is optional.}
    \label{tab:gpio-registers}
\end{table}

The modes are numbers, not types:

\begin{minted}{rust}
// Measured from the PAC's field encodings, and confirmed by the effect on the pin.
const DISABLED: u32 = 0;   // the pad is off; an analog peripheral may claim it
const INPUT:    u32 = 1;   // input buffer on, no pull
const PUSH_PULL:u32 = 4;   // driven high or low
const WIRED_AND: u32 = 8;  // open drain: 0 pulls low, 1 lets the pull-up win
const WIRED_OR:  u32 = 9;  // open drain, inverted sense
\end{minted}

Two of these decide most of the code in this book. \codeinline{PUSH_PULL} drives an LED and a
piezo. \codeinline{WIRED_AND} is what makes a bit-banged I2C bus possible, because ``output 1''
means ``release the line'' and the external pull-up makes it a real 1 — which is exactly what an
open-drain bus requires and exactly what a push-pull pin must never do.

\subsection{The Clock Gate Comes First}

The most expensive mistake in this whole project was not a wrong bit. It was touching a register
whose bus clock was off:

\begin{minted}{rust}
// Before: the write faulted, and the fault handler did not exist yet.
gpio.porta_dout().modify(|_, w| w.dout7().set_bit());

// After: gate the port, then talk to it.
dp.cmu_s.clken0().modify(|_r, w| w.gpio().set_bit());
gpio.porta_dout().modify(|_, w| w.dout7().set_bit());
\end{minted}

The measurement, taken from the fault status register after the fact:
\codeinline{CFSR = 0x8200} (a precise bus fault) with \codeinline{BFAR} pointing at the GPIO address
that was written. A GPIO register with no clock is not ``reads as zero'' — it is a bus fault, and if
the vector table still points at a default handler the board simply stops.

\subsection{Not Every Port Has Every Register}

\begin{minted}{rust}
// This package is a 48-pin QFN: PORTB and PORTD expose eight pins, not sixteen.
// MODEH covers 8..15, so on those two ports it does not exist at all.
match port {
    Port::A | Port::C => { /* MODEL and MODEH */ }
    Port::B | Port::D => { /* MODEL only */ }
}
\end{minted}

That is not a simplification invented for this book; it is what the SVD says, and a HAL that
pretends otherwise fails to compile \emph{after} someone has written the code. Two of the eight
compile errors during this project's pin primitive came from exactly that: the generated register
blocks differ per port, and a macro written for one does not typecheck for the other.

\section{A Pin Is Three Numbers}
\label{sec:pin-three-numbers}

The type-state approach to GPIO is seductive: make an output pin a different \emph{type} from an
input pin, make a pin claimed by one driver unclaimable by another, and let the compiler refuse the
mistakes. This book's repository had exactly that — 605 lines of it — and it is now retired, for a
reason worth stating precisely.

\subsection{What the Small Version Looks Like}

The whole abstraction is a port, a bit and a polarity:

\begin{minted}{rust}
pub enum Port { A, B, C, D }

pub struct Output<'a> {
    gpio: &'a pac::gpio_s::RegisterBlock,
    port: Port,
    bit: u8,
    active_low: bool,
}

impl<'a> Output<'a> {
    pub fn on(&self)  { self.write(!self.active_low) }
    pub fn off(&self) { self.write(self.active_low) }
    pub fn toggle(&self) {
        let current = read_dout(self.gpio, self.port);
        self.write(current & (1 << self.bit) == 0);
    }
}
\end{minted}

Everything else is a register write:

\begin{minted}{rust}
fn write_dout(gpio: &pac::gpio_s::RegisterBlock, port: Port, bit: u8, high: bool) {
    // Read-modify-write, always: DOUT is one register for sixteen pins.
    let set = |r: u32, w: &mut pac::W<_>| {
        let mask = 1u32 << bit;
        let value = if high { r | mask } else { r & !mask };
        unsafe { w.bits(value) }
    };
    match port {
        Port::A => gpio.porta_dout().modify(set),
        Port::B => gpio.portb_dout().modify(set),
        Port::C => gpio.portc_dout().modify(set),
        Port::D => gpio.portd_dout().modify(set),
    }
}
\end{minted}

That is the entire output path: 40 lines including the four-way match. The input path is shorter
still — \codeinline{DIN} needs no mask, only a shift — and the mode setter is the read-modify-write
of \codeinline{MODEL} or \codeinline{MODEH} shown in the previous section.

\subsection{Why the Larger Version Was Retired}

The type-state layer was not wrong in principle. It was wrong here, in three ways that only show up
in a real project:

\begin{enumerate}
    \item \textbf{Nothing used it.} Its only consumer was a board module that a newer facade
          replaced. An API with zero users compiles, lints clean, and rots: at the time it was
          retired, its README documented the user LED as \codeinline{PB2} active-\emph{high}, when
          the board's LED is \codeinline{PA7} active-\emph{low}.
    \item \textbf{It moved the wiring into the HAL.} To hand out a typed \codeinline{LED} the HAL had
          to know which pin the LED was on — a board fact in the wrong layer. The layered rule from
          section~\ref{sec:layer-boundaries} is that the HAL names no board wiring, and the typed
          pins were the mechanism by which it did.
    \item \textbf{It cost more than it prevented.} The three real defects in this project's GPIO
          code — a blind \codeinline{MODEL} write, a mode value with the wrong width, and an
          open-drain pin configured push-pull — were all invisible to the type system, because all
          three are legal code with the wrong \emph{number} in it. Types catch the mistakes you can
          name; registers catch the rest.
\end{enumerate}

\subsection{What Replaced the Type System}

An ownership rule that is checkable without generics: the board takes the peripherals once
(\codeinline{Board::take()}), and every pin comes out of it with a port and a bit named at the call
site:

\begin{minted}{rust}
let board = Board::take().expect("peripherals taken twice");
let led     = board.output_on(Port::A, 7, true);      // active low
let button  = board.input_on(Port::C, 1);             // pull-up fitted on the base
let buzzer  = board.output_on(Port::C, 3, false);
\end{minted}

The exclusivity that the type system used to provide is now the PAC singleton's job — one
\codeinline{take()} per program — and the wiring that the HAL used to hide lives in whichever crate
owns the board. Fewer lines, fewer layers, and the two facts a reader needs (which pin, which
polarity) are in the code they are reading.

\section{Read-Modify-Write, or Why the Bus Stopped Working}
\label{sec:gpio-rmw}

One register, sixteen pins. That sentence is the whole section, and it caused one of the most
confusing failures in this project.

\subsection{The Defect}

Setting a pin's mode looks like a write. It is a write to a register that also holds fifteen other
pins:

\begin{minted}{rust}
// Wrong: this configures *every* pin on the low half of PORTC, not just PC3.
gpio.portc_model().write(|w| w.mode3().pushpull());

// Right: read what is there, change one field, write it back.
gpio.portc_model().modify(|_r, w| w.mode3().pushpull());
\end{minted}

The wrong version does not fail. It reconfigures the fifteen other pins to their reset value — and on
this board the other pins on that half-port include the two that carry an I2C bus. A bus that had
just been proved working stops answering, with no error anywhere, because the pins that carried it
went back to \codeinline{DISABLED}.

\subsection{How It Was Found}

The failure looked like an I2C problem: the device stopped acknowledging, a scan returned nothing,
and reading the bus's own registers showed the peripheral configured exactly as intended. The
registers that mattered were the \emph{pin} ones:

\begin{minted}{text}
PORTC_MODEL = 0x00000000      # every pin on bits 0..7 disabled
PORTC_DOUT  = 0x00000030      # the bus lines still latched high
\end{minted}

Both lines at once: the latch said high, the pads were off. That combination can only come from a
write that reset the mode register, and a search for \codeinline{.write(} on that register found it
in one place — the function that was setting a buzzer pin.

\subsection{The Rule}

Every write to a GPIO configuration register in this project now goes through
\codeinline{.modify()}. The only \codeinline{.write()} calls are to registers that genuinely hold one
field and no other pin — and even then, the code says so in a comment, because the next person will
copy the pattern somewhere it does not hold.

The same rule applies to \codeinline{DOUT}: a blind \codeinline{write} there sets every pin on the
port high or low at once. It is the reason the pin primitives in this repository all take a port and
a bit, and never a whole register value:

\begin{minted}{rust}
// The signature is the rule: there is no way to write DOUT without saying which bit.
fn write_dout(gpio: &RegisterBlock, port: Port, bit: u8, high: bool)
\end{minted}

\subsection{A Second, Quieter Version of the Same Bug}

\codeinline{modify()} on a register whose clock is off is a bus fault, but \codeinline{modify()} on
a register with a \emph{write-only field} is worse: it reads back a value the field does not
contribute to, so the read-modify-write silently clears it. This is why the mode values in this
project are written as an explicit numeric field of the right width rather than as a chain of
boolean writers, and why the IADC's configuration registers (chapter 7's sibling problem) are written
once, in one function, in a documented order — conditional writes into a synchronising peripheral
lose the field you did not intend to touch.

\section{Open Drain, and the Line That Has Two Masters}
\label{sec:gpio-open-drain}

An I2C bus has two wires and, on each of them, two possible masters, each of which may want to pull
the line low while the other wants it high. The electrical answer is the reason the mode list has a
\codeinline{WIRED\_AND} in it.

\subsection{What WIRED\_AND Does}

\begin{minted}{text}
mode = WIRED_AND   (8)
  DOUT bit = 0  ->  the pad pulls the line to ground
  DOUT bit = 1  ->  the pad releases the line (high impedance)
\end{minted}

The line is held high by an external resistor. So ``1'' is not something the pin produces — it is
something the pin \emph{stops preventing}. Two devices can then share the wire: any one of them can
pull it low, and the only way it goes high is if nobody does. That is what makes arbitration
possible, and it is why a push-pull pin must never be connected to a shared bus: driving it high
against another device driving low is a short circuit through two transistors.

\subsection{The Measurement That Proved the Pull-Ups}

The IMU could not be found for two sessions. Before blaming the pin map, this project measured
whether the bus lines had pull-ups at all:

\begin{minted}{text}
# every pin on PORTB set to input, then read DIN
PORTB_MODEL = 0x11111111     # all inputs
PORTB_DIN   = 0x0000000C     # bits 2 and 3 high
\end{minted}

Bits 2 and 3 are \codeinline{PB2} and \codeinline{PB3}: the two pins of the IMU's bus, reading high
with nothing driving them. That is an external pull-up, present and working — which retired the
``the board has no pull-ups'' hypothesis in one instruction and moved the investigation to the pin
\emph{order}.

\subsection{Bit-Banging an Open-Drain Bus}

With \codeinline{WIRED\_AND} the whole I2C protocol reduces to two writes and one read:

\begin{minted}{rust}
fn sda(&self, high: bool) { write_dout(self.gpio, self.bus.port, self.bus.sda_bit, high) }
fn scl(&self, high: bool) { write_dout(self.gpio, self.bus.port, self.bus.scl_bit, high) }
fn read_sda(&self) -> bool { read_din(self.gpio, self.bus.port) & (1 << self.bus.sda_bit) != 0 }
\end{minted}

\textbf{There is no direction switch.} On an open-drain bus a pin used as an input with the output
latch high is the same electrical state as a released output — which is why a bit-banged I2C master
can be four functions and no state machine, and why the same pin can be written and read in the same
transfer.

\subsection{Half a Bit Period, and Why It Is a Number in Cycles}

Timing on a bit-banged bus is a delay per half period, and the delay is not a duration — it is a
count of loop iterations whose real length depends on the compiler:

\begin{minted}{rust}
// ~5 us at 19 MHz, measured: 60 iterations of this loop on this compiler.
const HALF_BIT_CYCLES: u32 = 60;
\end{minted}

That constant is a measurement, not a datasheet value, and it is the reason chapter~14 exists: the
same class of mistake — \codeinline{asm::delay(n)} assumed to be \codeinline{n} cycles when it is
about three per iteration — turned a 2.7 kHz tone into something else entirely.

\section{What the Vendor's Own Documents Get Wrong}
\label{sec:gpio-vendor-errors}

Three documents describe this board's pins, and two of them disagree with the board. This section is
short and uncomfortable, and it exists because a wrong pin in a comment is worse than no comment: it
is believed.

\subsection{The User LED}

\begin{table}[htbp]
    \centering
    \footnotesize
    \begin{tabular}{@{}p{0.34\textwidth}p{0.26\textwidth}p{0.32\textwidth}@{}}
        \toprule
        \textbf{Source} & \textbf{Says} & \textbf{What it is} \\
        \midrule
        The HAL's retired GPIO README & \codeinline{PB2}, active high & wrong pin, wrong polarity \\
        An earlier board module doc   & \codeinline{PB02/D10}       & wrong pin \\
        Register read-back (measured) & \codeinline{PA7}, active low & \textbf{this one} \\
        \bottomrule
    \end{tabular}
    \caption{The user LED, as documented and as measured. The measurement is a write to
    PORTA\_DOUT bit 7 and a read of PORTA\_DOUT to confirm the latch moved.}
    \label{tab:led-truth}
\end{table}

Why the error survived: nothing in the project ever \emph{used} the wrong pin in a way that failed.
A pin that does not exist toggles a latch nobody looks at, and the code still returns
\codeinline{Ok}. Only the blink example — the one that must be seen — exposed it.

\subsection{The IMU's Two Wires}

The Arduino variant for this board declares:

\begin{minted}{text}
PIN_WIRE_SDA1 = PB3
PIN_WIRE_SCL1 = PB2
\end{minted}

The board is wired the other way round. Measured, in this order:

\begin{enumerate}
    \item a bit-banged scan of all 128 addresses on the pair \codeinline{(PB2, PB3)} with
          SDA=\codeinline{PB2} found a device at \codeinline{0x6A};
    \item the same scan with the roles swapped found nothing;
    \item the device at \codeinline{0x6A} answered \codeinline{WHO\_AM\_I = 0x6A}, which is the
          LSM6DS3TR-C's identity register.
\end{enumerate}

Three independent facts — the scan, the negative control, and the identity register — agree. The
variant file does not, and it is the file a reader would trust first.

\subsection{The Lesson, Which Is Not ``the Vendor Is Wrong''}

The vendor's documents are mostly right and are the only reason this works at all: the SVD, the
reference manual, the bootloader layout and the radio headers are all correct and all essential. But
the pin map of a \emph{board} is not the same document as the pin map of a \emph{chip}, and it is
maintained by people assembling a product, not by the silicon team.

So the rule this project adopted: \textbf{every pin fact is measured on the board, and the
measurement is written next to the claim}. Two of the three corrections above cost an afternoon
each. Both would have been found in ten minutes by reading a register instead of a comment.

\section{Proving a Pin Without Trusting Your Eyes}
\label{sec:gpio-proving}

The blink is the first program every embedded developer writes, and it is the first one that lies. A
blink that does not blink has four different causes that look identical: the code never ran, the pin
is wrong, the polarity is wrong, or the LED is not connected. This section is about telling them
apart with instruments that are already there.

\subsection{The Case That Made This Necessary}

This project had a blink that flashed successfully and did not blink. The write verified, the
read-back matched, the board was silent. Every step of the toolchain reported success:

\begin{minted}{text}
cargo build            Finished `release` profile [optimized] target(s)
python3 scripts/elf-load-range.py   load_min=0x08006000 load_max=0x08006813
probe-rs download      Finished in 0.41s
probe-rs verify        Finished in 0.28s
\end{minted}

Four green steps and no light. The only way forward was to stop asking the board to show and start
asking the register to tell.

\subsection{Reading the Latch Over SWD}

The probe can read any address while the program runs. The output latch is one word, and it toggles
or it does not:

\begin{minted}{bash}
# halt, read DOUT twice, resume
openocd -c "init; halt" -c "mdw 0x5003C040 1" -c "sleep 500" -c "mdw 0x5003C040 1" -c "resume; shutdown"

# 0x5003c040: 00000080      <- bit 7 high
# 0x5003c040: 00000000      <- bit 7 low, 500 ms later
\end{minted}

That pair of readings is a proof of the \emph{program}: the latch is moving at the rate the code
intended. What it is not is a proof of the \emph{board}: a latch that toggles and an LED that stays
dark means the pin is not the LED's pin, or the LED is active-low and the code drives it backwards.

The second reading of \codeinline{PORTA\_MODEL} closes the other half:

\begin{minted}{text}
PORTA_MODEL = 0x00000040   # mode 4 (push-pull) in the field for pin 7
\end{minted}

Latch moving, mode correct, and still no light: that combination is now a board question, not a
software one — which is exactly the division of labour a measurement is for.

\subsection{Three Instruments, Three Questions}

\begin{table}[htbp]
    \centering
    \footnotesize
    \begin{tabular}{@{}p{0.26\textwidth}p{0.34\textwidth}p{0.34\textwidth}@{}}
        \toprule
        \textbf{Instrument} & \textbf{Question it answers} & \textbf{How it is read} \\
        \midrule
        The LED & is the whole chain alive? & eyes, and only after the two below \\
        \codeinline{DOUT}/\codeinline{MODEL} over SWD & is the program right? & \codeinline{mdw} on the port \\
        A mailbox in RAM & did the code reach this point? & \codeinline{mdw} on a magic word \\
        \bottomrule
    \end{tabular}
    \caption{Order matters: the LED answers last, because it answers slowly and ambiguously.}
    \label{tab:gpio-instruments}
\end{table}

\subsection{The Mailbox, Which Is the Real Lesson}

The instrument that made the rest of this project possible is not a pin at all. It is a small block
of RAM — a magic word, a sequence number, a status and a message — that a program writes before it
does anything it might not survive:

\begin{minted}{rust}
// 0x2003F000: magic "MBOX", seq, status (0 info / 1 pass / 2 fail), len, message bytes
mailbox::write(Status::Pass, b"cap=0x20000000 4000@400Hz real=clean");
\end{minted}

Read with one \codeinline{mdw} over SWD, it answers questions that no LED can: did the classifier
load its calibration, how many samples did the ring hold, which of the two I2C aliases answered. It
costs 300 bytes of RAM and it replaced an entire class of guesswork — which is why every example in
this book reports through it rather than through a serial port that may not exist.


\clearpage

\begin{keytakeaways}
    \item A GPIO port on this part is five registers: \codeinline{MODEL}, \codeinline{MODEH},
          \codeinline{DOUT}, \codeinline{DIN} and the port's clock gate.
    \item \codeinline{MODEH} does not exist on every port: on PORTB and PORTD this package has eight
          pins, so those pins live in \codeinline{MODEL} alone.
    \item A pin is a port, a bit and a polarity. That is enough to drive an LED, read a button and
          open-drain a bus, and it is 189 lines — not 605.
    \item Writing \codeinline{MODEL} or \codeinline{MODEH} without reading first reconfigures every
          other pin on the same half-port, which is how a working bus stops working.
    \item The LED on this board is \codeinline{PA7}, active \emph{low}. Two of the vendor's own
          documents disagree with that.
    \item A blink is proven by the register toggling, not by the LED looking on.
\end{keytakeaways}
